\documentclass[12pt]{article}
\usepackage{amsmath,amssymb}
\usepackage[margin=1in]{geometry}
\usepackage{setspace}
\onehalfspacing
\date{}

% ---------------------------------------------------------------------
% CHANGES FROM PREVIOUS VERSION
% 1. Single closed economy: country index i dropped everywhere.
% 2. Final good is the numeraire (P_t = 1). Developer utility and mu are
%    therefore in units of the final good without a P divisor.
% 3. Consistency fixes found in the market-clearing review:
%    a. Household profit income was missing the R&D wage bill paid by
%       innovators; without it Walras' law fails. Now included.
%    b. Final-good market clearing added: Y = C + fuel. This closes the
%       system and is implied by the other conditions (Walras).
%    c. Interest rate r was free; now pinned by the household Euler
%       equation (firms discount with the household's rate).
%    d. The quality state is the (sigma-1)-moment aggregate Q^z, not the
%       mean q-bar: profits and the price index depend on q^(sigma-1).
%       Random targeting makes Q^z a sufficient statistic with an exact
%       law of motion.
%    e. Patent value made exact: V(q) = q^(sigma-1) v, with the innovator's
%       advantage decaying as the sector aggregate grows. Innovation and
%       displacement timing aligned (research at t, productive at t+1).
%    f. Dispatch: interior case (all plants run) stated as an assumption.
% 4. New Section 8: market clearing and equilibrium definition.
% ---------------------------------------------------------------------

\begin{document}

\section{Environment}

Time is discrete, $t=0,1,2,\dots$. The economy is closed. Population
$L_t$ is exogenous and supplied inelastically. The final good is the
numeraire, $P_t=1$. The final good is produced using labour and
electricity. Electricity is produced by installed power plants of
different types $k\in\mathcal K$ (for example solar, wind, gas, coal),
partitioned into a clean set $\mathcal C$ and a fossil set $\mathcal F$.
Each period a mass of developers operates; each developer chooses which
type of plant to build, or not to build, according to a random-utility
model over project values. Building a clean plant requires a continuum of
differentiated clean capital varieties (panels, turbines, inverters);
building a fossil plant requires a separate continuum of dirty capital
varieties (boilers, turbines, gas trains), and operating it requires
fuel. Each variety firm produces one unique variety. Innovation raises
the quality of a variety, lowering its price, lowering the capital cost
of the corresponding plant type, and raising demand for the variety.

Labour is used in final-good production, variety production and
research:
\[
L_t = L^Y_t + L^M_t + S_t,
\qquad S_t=S^c_t+S^d_t,
\]
where $S^z_t$ is the mass of scientists working on sector
$z\in\{c,d\}$ varieties.

\section{Households}

The representative household owns all variety firms, plants, development
projects and research firms, and chooses consumption $\{C_t\}$ to
maximise
\[
\max_{\{C_t\}}\sum_{t=0}^{\infty}\frac{1}{(1+\rho)^t}\,u(C_t)
\qquad\text{s.t.}\qquad
C_t + T_t = w_tL_t + \Pi_t ,
\]
where $T_t$ is a lump-sum tax and $\Pi_t$ is the aggregate net cash flow
of all firms owned by the household:
\begin{align*}
\Pi_t={}&
\underbrace{\sum_{z\in\{c,d\}}\int_0^{M^z}\pi^z_t(\nu)\,d\nu}
_{\text{variety-firm profits}}
+\underbrace{\sum_{k\in\mathcal C}A_{kt}K_{kt}P^F_{kt}
+\sum_{k\in\mathcal F}A_{kt}K_{kt}\big(P^E_t-c^k_t\big)}
_{\text{plant operating profits}}\\
&-\underbrace{\sum_{k\in\mathcal C\cup\mathcal F}F^k_t\,s_{kt}D_t}
_{\text{investment outlays}}
-\underbrace{\big[(w_t-\xi_t)S^c_t+w_tS^d_t\big]}
_{\text{R\&D outlays net of subsidy}} .
\end{align*}
There is no traded asset, so the Euler equation defines the discount rate
that firms use:
\[
1+r_{t+1}=(1+\rho)\,\frac{u'(C_t)}{u'(C_{t+1})} .
\]

\section{Final Good}

Competitive producers choose labour $L^Y_t$ and electricity $E_t$ to
maximise profit:
\[
\max_{L^Y_t,\,E_t}\quad Y_t - w_tL^Y_t - P^E_tE_t
\qquad\text{s.t.}\qquad
Y_t = \left(L^Y_t\right)^{1-\beta}E_t^{\,\beta}
\]
with first-order conditions $w_tL^Y_t=(1-\beta)Y_t$ and
$P^E_tE_t=\beta Y_t$.

\section{Power Plant Development}\label{sec:dev}

\subsection{Plants and electricity supply}

A unit of installed capacity of type $k$ produces $A_{kt}$ units of
electricity per period, where $A_{kt}$ is the exogenous capacity factor
(irradiance, wind, availability). Capacity depreciates at rate
$\delta_k$. Clean plants have no variable cost. Fossil plants burn fuel
$k$, supplied by a domestic resource sector at constant marginal cost
$p^k_t$ units of the final good per unit of fuel, converted at thermal
efficiency $h_k$ and taxed at the carbon tax $\tau^E_t$ scaled by
emissions intensity $e_k$:
\[
c^k_t=0 \quad (k\in\mathcal C),
\qquad
c^k_t=\frac{(1+\tau^E_te_k)\,p^k_t}{h_k}
\quad (k\in\mathcal F).
\]
Electricity is homogeneous. Fossil plants sell at the competitive market
price $P^E_t$ and run whenever $P^E_t\ge c^k_t$. Clean plants of type
$k$ receive a feed-in tariff $P^F_{kt}$ per unit of output, with
$P^F_{kt}=P^E_t$ when no support is in place, and always run. We restrict
attention to the interior case $P^E_t\ge c^k_t$ for all
$k\in\mathcal F$, so all installed capacity runs:
\[
E_t=\sum_{k\in\mathcal K}A_{kt}K_{kt}.
\]

\subsection{Capital cost of a plant}

Building one unit of capacity of type $k$ requires $\iota_k$ units of a
CES bundle of capital varieties from the corresponding sector, clean
($z=c$) for $k\in\mathcal C$ and dirty ($z=d$) for $k\in\mathcal F$:
\[
X^z_t=\left[\int_0^{M^z}
x^z_t(\nu)^{(\sigma-1)/\sigma}\,d\nu\right]^{\sigma/(\sigma-1)},
\qquad z\in\{c,d\},
\]
with price index $P^z_t$ defined in Section~\ref{sec:varieties}.
Capital cost per unit of capacity is
\[
F^k_t=\iota_k\,P^c_t\quad(k\in\mathcal C),
\qquad
F^k_t=\iota_k\,P^d_t\quad(k\in\mathcal F).
\]

\subsection{Value of a project}

Let $Q^k_t$ be the present value of operating profits from one unit of
type-$k$ capacity that is in place at $t$:
\[
Q^k_t=A_{kt}\big(P^F_{kt}-c^k_t\big)
+\frac{1-\delta_k}{1+r_{t+1}}\,Q^k_{t+1}
\quad(k\in\mathcal C),
\qquad
Q^k_t=A_{kt}\big(P^E_t-c^k_t\big)
+\frac{1-\delta_k}{1+r_{t+1}}\,Q^k_{t+1}
\quad(k\in\mathcal F).
\]
A plant built at $t$ becomes operational at $t+1$. The value of
developing one unit of type-$k$ capacity at $t$ is
\[
W^k_t=-F^k_t+\frac{1}{1+r_{t+1}}\,Q^k_{t+1},
\qquad k\in\mathcal C\cup\mathcal F,
\]
and the value of not building (the outside option, $k=0$) is $W^0_t=0$.

\subsection{Developer random utility}

Each period a mass $D_t$ of developers $\omega$ operates, each holding
one unit-capacity project opportunity. Developer $\omega$'s payoff from
choosing type $k\in\mathcal K_0\equiv\{0\}\cup\mathcal C\cup\mathcal F$
is
\[
U_{kt}(\omega)=W^k_t+\mu\,\epsilon_{k}(\omega),
\]
where $\epsilon_{k}(\omega)$ is an idiosyncratic project-specific
component (site quality, land, permitting, grid connection, financing)
and $\mu>0$ scales its dispersion in units of the final good. The
developer chooses $k$ to maximise $U_{kt}(\omega)$.

\subsection{Nested logit shares}

The shocks $\epsilon_{k}(\omega)$ follow a generalised extreme value
distribution with nests $g\in\{\{0\},\mathcal C,\mathcal F\}$ and nesting
parameters $\lambda_g\in(0,1]$. Writing $v^k_t\equiv W^k_t/\mu$, the
share of developers choosing type $k$ in nest $g$ is
\[
s_{kt}=s_{k|g,t}\cdot s_{gt},
\qquad
s_{k|g,t}=\frac{\exp(v^k_t/\lambda_g)}
{\sum_{m\in g}\exp(v^m_t/\lambda_g)},
\qquad
s_{gt}=\frac{\exp(I_{gt})}{\sum_{h}\exp(I_{ht})},
\]
with nest inclusive values
\[
I_{gt}=\lambda_g\ln\sum_{m\in g}\exp\big(v^m_t/\lambda_g\big).
\]
When $\lambda_g=1$ for all $g$ this is the multinomial logit
\[
s_{kt}=\frac{\exp(v^k_t)}{1+\sum_{m\in\mathcal C\cup\mathcal F}
\exp(v^m_t)},
\qquad
\frac{\partial\ln s_{kt}}{\partial v^m_t}=\mathbf 1\{k=m\}-s_{mt}.
\]
Aggregate new capacity of type $k$ is $s_{kt}D_t$, and the clean and
fossil build shares are $s^{\mathcal C}_t=\sum_{k\in\mathcal C}s_{kt}$
and $s^{\mathcal F}_t=\sum_{k\in\mathcal F}s_{kt}$.

\subsection{Capacity accumulation}

\[
K_{k,t+1}=(1-\delta_k)K_{kt}+s_{kt}D_t,
\qquad k\in\mathcal C\cup\mathcal F .
\]
Capacity is a state variable: developers' current choices determine
future electricity supply, and the forward-looking project value
$Q^k_{t+1}$ depends on future capacity through $P^E$.

\subsection{Demand for capital varieties}

Expenditure on each variety sector is the corresponding capex flow:
\[
R^c_t\equiv P^c_tX^c_t=\sum_{k\in\mathcal C}\iota_k\,P^c_t\,s_{kt}D_t,
\qquad
R^d_t\equiv P^d_tX^d_t=\sum_{k\in\mathcal F}\iota_k\,P^d_t\,s_{kt}D_t.
\]
Developers minimise the cost of each bundle, so demand for variety $\nu$
of sector $z$ is
\[
x^z_t(\nu)=\left(\frac{p^z_t(\nu)}{P^z_t}\right)^{-\sigma}
\frac{R^z_t}{P^z_t} .
\]
This is the market-size channel of directed innovation: lower $P^c_t$
lowers $F^k_t$ for all $k\in\mathcal C$, raises $W^k_t$, raises the clean
build share $s^{\mathcal C}_t$, and raises the revenue $R^c_t$ of clean
variety firms relative to $R^d_t$.

\subsection{Policy channels}

A carbon tax raises $c^k_t$ for $k\in\mathcal F$, lowers $Q^k$ and $W^k$
for fossil plants, and shifts developers toward clean types and the
outside option at a rate governed by $\mu$. A feed-in tariff
$P^F_{kt}>P^E_t$ raises $Q^k$ and $W^k$ for clean plants directly. Both
raise $R^c_t$ relative to $R^d_t$ and therefore redirect innovation
toward clean varieties through market size; the R\&D subsidy $\xi_t$
(Section~\ref{sec:innovation}) acts on the free-entry condition of
scientists instead. Within the nested model, policies that change
relative project values inside the clean nest (a solar-specific feed-in
tariff) reallocate developers within the nest at rate
$1/\lambda_{\mathcal C}$ and across nests at rate $1$.

\paragraph{Remark (free entry of developers).}
An alternative to an exogenous $D_t$ is free entry: a developer pays a
development cost $\kappa^D_t$ to draw a project, and entry continues
until the expected value of a draw, $\mu\ln\sum_g\exp(I_{gt})$, equals
$\kappa^D_t$.

\paragraph{Remark (estimation).}
With the outside option as reference, Berry's inversion gives
\[
\ln s_{kt}-\ln s_{0t}
=\frac{W^k_t}{\mu}
\quad\Big(\text{nested: }+\,(1-\lambda_g)\ln s_{k|g,t}\Big),
\]
so $\mu$ and $\lambda_g$ are identified from plant-level data on
commissioned projects (shares of new capacity by type) together with
observed capital costs and expected operating margins, with unobserved
type-time project value absorbed by fixed effects.

\section{Capital Variety Firms}\label{sec:varieties}

\subsection{Firm structure}

In each sector $z\in\{c,d\}$ there is a fixed mass $M^z$ of active
firms, each producing one unique variety $\nu$ under patent with current
quality $q^z_t(\nu)$. Quality reflects manufacturing efficiency: higher
quality means the firm can produce more units per worker.

\subsection{Production and pricing}

Firm $\nu$ in sector $z$ produces using $\psi/q^z_t(\nu)$ units of labour
per unit of output, so marginal cost is
\[
mc^z_t(\nu) = \frac{\psi\,w_t}{q^z_t(\nu)}
\]
and the firm sets a constant markup over marginal cost:
\[
p^z_t(\nu) = \frac{\sigma}{\sigma-1}\cdot\frac{\psi\,w_t}{q^z_t(\nu)} .
\]

\subsection{Quality aggregate and price index}

Define the sector quality aggregate
\[
\mathcal Q^z_t\equiv\int_0^{M^z} q^z_t(\nu)^{\sigma-1}\,d\nu .
\]
The price index is
\[
P^z_t = \left[\int_0^{M^z}
p^z_t(\nu)^{1-\sigma}\,d\nu\right]^{1/(1-\sigma)}
=\frac{\sigma}{\sigma-1}\cdot\psi\,w_t\cdot
\big(\mathcal Q^z_t\big)^{1/(1-\sigma)} .
\]
Higher $\mathcal Q^c_t$ lowers clean capital costs $F^k_t$,
$k\in\mathcal C$, and raises the clean build share
$s^{\mathcal C}_t$; symmetrically for $\mathcal Q^d_t$ and fossil
plants.

\subsection{Firm profits and labour demand}

Flow profit of a firm with quality $q$ in sector $z$:
\[
\pi^z_t(q) = \frac{1}{\sigma}
\left(\frac{p^z_t(q)}{P^z_t}\right)^{1-\sigma}R^z_t
=\frac{q^{\sigma-1}}{\mathcal Q^z_t}\cdot\frac{R^z_t}{\sigma} .
\]
Sector profits sum to $R^z_t/\sigma$ and the sector wage bill is
$(\sigma-1)R^z_t/\sigma$, so labour used in variety production is
\[
L^M_t=\sum_{z\in\{c,d\}}\int_0^{M^z}\frac{\psi\,x^z_t(\nu)}{q^z_t(\nu)}\,d\nu
=\frac{\sigma-1}{\sigma}\cdot\frac{R^c_t+R^d_t}{w_t} .
\]

\section{Innovation}\label{sec:innovation}

\subsection{Quality ladder and timing}

A scientist working in period $t$ targets a randomly drawn firm $\nu$ in
sector $z$ and succeeds with probability $\eta$. A successful innovator
raises the target's quality to $(1+\gamma)q^z_t(\nu)$, displaces the
incumbent at the end of period $t$, and produces with the patent from
$t+1$ onward. With $S^z_t$ scientists in sector $z$, each incumbent is
displaced with probability $\eta S^z_t/M^z$ (assumed below $1$). By
random targeting, the expected contribution of a success to the quality
aggregate is proportional to the average of $q^{\sigma-1}$, so
\[
\mathcal Q^z_{t+1}=\mathcal Q^z_t\left[1+\frac{\eta S^z_t}{M^z}
\Big((1+\gamma)^{\sigma-1}-1\Big)\right] .
\]

\subsection{Value of a patent}

The value at $t$ of a patent with quality $q$ in sector $z$ is
\[
V^z_t(q)=\pi^z_t(q)+\frac{1-\eta S^z_t/M^z}{1+r_{t+1}}\,V^z_{t+1}(q) .
\]
Since $\pi^z_t(q)$ is linear in $q^{\sigma-1}$, so is $V^z_t(q)$:
$V^z_t(q)=q^{\sigma-1}v^z_t$ with
\[
v^z_t=\frac{R^z_t}{\sigma\,\mathcal Q^z_t}
+\frac{1-\eta S^z_t/M^z}{1+r_{t+1}}\,v^z_{t+1} .
\]
A patent's relative profit decays as the sector aggregate
$\mathcal Q^z$ grows around it.

\subsection{Scientist optimisation}

Workers are freely mobile between production and research, so
scientists earn the wage $w_t$ and are hired by competitive research
firms. The government pays a per-scientist clean R\&D subsidy $\xi_t$
to research firms for each scientist targeting clean varieties. A
scientist who succeeds at $t$ holds, from $t+1$, a patent of quality
$(1+\gamma)q^z_t(\nu)$ for a randomly drawn $\nu$, with expected
$q^{\sigma-1}$ equal to $\mathcal Q^z_t/M^z$. Free entry into each
sector implies
\[
\eta\,(1+\gamma)^{\sigma-1}\,\frac{\mathcal Q^c_t}{M^c}\cdot
\frac{v^c_{t+1}}{1+r_{t+1}}+\xi_t\le w_t,
\qquad
\eta\,(1+\gamma)^{\sigma-1}\,\frac{\mathcal Q^d_t}{M^d}\cdot
\frac{v^d_{t+1}}{1+r_{t+1}}\le w_t,
\]
each holding with equality whenever $S^z_t>0$. When both hold with
equality, scientists allocate across sectors until the subsidised
expected value of a clean innovation equals that of a dirty one.

\subsection{Innovation incentive}

Along a path where all firms in a sector have the same quality, the
expected flow profit from a successful innovation in sector $z$ is
\[
\Pi^z_t=\eta\,(1+\gamma)^{\sigma-1}\cdot\frac{R^z_t}{\sigma M^z},
\qquad
\frac{\Pi^c_t}{\Pi^d_t}=\frac{R^c_t/M^c}{R^d_t/M^d},
\]
the ratio of capex revenue per variety across the two sectors. Carbon
taxes and feed-in tariffs raise $R^c_t/R^d_t$ by shifting developers
toward clean plants and operate on innovation through this market-size
ratio, with strength governed by $\mu$; the R\&D subsidy $\xi_t$
operates directly on the free-entry condition.

\section{Government}

The government levies a carbon tax on fuel and a lump-sum tax, and pays
the feed-in tariff and the clean R\&D subsidy. Carbon-tax revenue is
\[
R^E_t=\sum_{k\in\mathcal F}\tau^E_te_k\,p^k_t\,\frac{A_{kt}K_{kt}}{h_k},
\]
and feed-in-tariff outlays are
\[
G^F_t=\sum_{k\in\mathcal C}\big(P^F_{kt}-P^E_t\big)A_{kt}K_{kt}.
\]
The budget balances each period:
\[
R^E_t+T_t=G^F_t+\xi_t\,S^c_t ,
\]
so that $T_t<0$ is a lump-sum rebate when revenue exceeds outlays.

\section{Market Clearing and Equilibrium}\label{sec:eq}

\subsection{Market clearing}

\begin{itemize}
\item Electricity: $P^E_tE_t=\beta Y_t$ with $E_t=\sum_kA_{kt}K_{kt}$,
so $P^E_t=\beta\,(L^Y_t)^{1-\beta}E_t^{\beta-1}$.
\item Labour: $L_t=L^Y_t+L^M_t+S^c_t+S^d_t$ with
$w_tL^Y_t=(1-\beta)Y_t$ and $L^M_t=(\sigma-1)(R^c_t+R^d_t)/(\sigma w_t)$.
\item Varieties: each firm's output equals demand $x^z_t(\nu)$; summing
$p^z_t(\nu)x^z_t(\nu)$ over $\nu$ gives $R^z_t$ by construction of the
CES bundle.
\item Fuel: fuel use is $A_{kt}K_{kt}/h_k$ for $k\in\mathcal F$,
produced from $p^k_t$ units of the final good per unit.
\item Final good: $Y_t=C_t+\sum_{k\in\mathcal F}p^k_t\,A_{kt}K_{kt}/h_k$.
\end{itemize}
Substituting the household budget, the definition of $\Pi_t$, the
final-good first-order conditions and the government budget gives
$C_t=Y_t-\sum_{k\in\mathcal F}p^k_tA_{kt}K_{kt}/h_k$, so the final-good
market clears by Walras' law. Capital varieties are produced from
labour, capex is a payment from developers to variety firms, and the
carbon tax, feed-in tariff, R\&D subsidy and lump-sum tax are transfers;
the only use of the final good other than consumption is fuel.

\subsection{Equilibrium definition}

Given initial capacities $\{K_{k0}\}$ and quality aggregates
$\{\mathcal Q^z_0\}$, exogenous paths $\{L_t,D_t,A_{kt},p^k_t\}$ and
policies $\{\tau^E_t,P^F_{kt},\xi_t\}$, an equilibrium is a sequence of
prices $\{w_t,P^E_t,P^c_t,P^d_t,r_{t+1}\}$, quantities
$\{C_t,Y_t,E_t,L^Y_t,L^M_t,S^c_t,S^d_t,R^c_t,R^d_t\}$, project values
$\{Q^k_t,W^k_t,s_{kt}\}$, patent values $\{v^z_t\}$, taxes $\{T_t\}$ and
states $\{K_{k,t+1},\mathcal Q^z_{t+1}\}$ such that: households satisfy
the Euler equation and budget; final-good producers satisfy their
first-order conditions; developers choose by the nested logit given
$\{W^k_t\}$; $Q^k_t$, $W^k_t$ and $v^z_t$ satisfy their recursions;
variety firms price at the markup and the price indices hold; scientists
satisfy free entry in each sector; the government budget balances; the
state variables follow their laws of motion; and the electricity, labour
and final-good markets clear.


\end{document}
